\chapter{منطق ریاضی}
\thispagestyle{empty}

\begin{quote}
	{\lotus
		هرگز در منطق، غافلگیر شدن نمی‌تواند وجود داشته باشد. \quad - ویتگنشتاین%
		\LTRfootnote{ Ludwig Josef Johann Wittgenstein}
	}
\end{quote}

\section{معماهای منطقی}
\begin{problem} \label{loprpedar}
	پدری به دو فرزندش (یک پسر و یک دختر) می‌گوید که در حیاط خانه بازی کنند ولی کثیف نشوند. با این حال، در حین بازی هر دو کودک روی پیشانی‌شان گل می‌مالند. وقتی بچه‌ها بازی را متوقف می‌کنند، پدر می‌گوید: «دست‌کم یکی از شما پیشانی‌اش گلی است.» سپس از بچه‌ها می‌خواهد که به این پرسش پاسخ «بله» یا «خیر» بدهند: «آیا می‌دانی پیشانی‌ات گلی است یا نه؟» پدر این پرسش را دو بار مطرح می‌کند. بچه‌ها در هر نوبت چه پاسخی خواهند داد؟ فرض کنید هر کودک می‌تواند ببیند که پیشانی خواهر یا برادرش گلی است یا نه، ولی نمی‌تواند پیشانی خود را ببیند. همچنین فرض کنید هر دو کودک راستگو هستند و پاسخ‌ها را هم‌زمان می‌دهند.
\end{problem}

\begin{problem} \label{loprabccolor}
	آقایان
	\lr{A}، \lr{B} و \lr{C}
	هر سه در منطق مهارت کامل دارند. هر یک می‌تواند کلیه‌ی نتیجه‌هایی را که از مقدماتی معین به دست می‌آید، بی‌درنگ استخراج کند. همچنین هر یک می‌داند که دو نفر دیگر نیز مانند خودش در منطق مهارت دارند. به هر یک از این آقایان دو تمبر قرمز، دو تمبر زرد و سه تمبر سبز نشان می‌دهیم. پس از آن، چشم‌های هر سه را می‌بندیم و به پیشانی هر کدام یک تمبر می‌چسبانیم. چهار تمبر باقیمانده را پنهان می‌کنیم. سپس چشم‌های این سه نفر را باز می‌کنیم. از آقای
	\lr{A}
	می‌پرسیم که آیا می‌تواند بگوید که تمبری که به پیشانی او چسبیده شده است چه رنگی را به طور قطع نخواهد داشت. او پاسخ منفی می‌دهد. همین پرسش را از آقای
	\lr{B}
	می‌کنیم. او نیز پاسخ منفی می‌دهد.
	آیا با استفاده از این اطلاعات می‌توان رنگ تمبری را که به پیشانی
	\lr{A} یا \lr{B} یا \lr{C}
	چسبیده شده است، استنتاج کرد؟
\end{problem}

\begin{problem} \label{loprmcgregor}
	آقای مک‌گرگور%
	\LTRfootnote{ McGregor}%
	، مأمور سرشماری، زمانی در جزیره‌ی شوالیه‌ها (همیشه راستگویان) و سرباز‌ها (همیشه دروغگویان)%
	\LTRfootnote{ Island of Knights and Knaves}
	تحقیق میدانی انجام داد. او تصمیم گرفت فقط با زوج‌های متأهل مصاحبه کند. \\
	مک‌گرگور در خانه‌ی اول را زد و مردی در را تا نیمه باز کرد و پرسید چه کار دارد. مک‌گرگور پاسخ داد: «من مأمور سرشماری هستم و درباره‌ی شما و همسرتان اطلاعات نیاز دارم. کدام یک از شما شوالیه و کدام یک از شما سرباز است؟» مرد با خشم جواب داد: «ما هر دو سرباز هستیم.» \\
	در خانه‌ی دوم، مک‌گرگور از مرد پرسید: «آیا هر دو سرباز هستید؟» شوهر پاسخ داد: «حداقل یکی از ما سرباز است.» \\
	در خانه‌ی سوم، مردی خجالتی در را باز کرد و پس از پرسش مک‌گرگور پاسخ داد: «اگر من شوالیه باشم، پس همسرم هم شوالیه است.» \\
	وقتی مک‌گرگور از چهارمین زوج سؤال پرسید، مرد گفت: «من و همسرم مثل هم هستیم، یا هر دو شوالیه‌ایم یا هر دو سربازیم.» \\
	در هر خانه مرد و زن از چه نوعی (شوالیه یا سرباز) هستند؟
\end{problem}

\begin{problem} \label{loprladytiger}
	پادشاه به زندانی توضیح داد که دو اتاق وجود دارد که در هر یک از این دو اتاق دختری زیبا یا ببری گرسنه است. اما ممکن است در هر دو اتاق ببر یا در هر دو اتاق دختر یا در یک اتاق ببر و در اتاق  دیگر دختر باشد.
	زندانی پرسید: «اگر در هر دو اتاق ببر باشد چطور؟ در آن صورت چه کار کنم؟»
	پادشاه پاسخ داد: «این از بخت بد شما خواهد بود.»
	زندانی پرسید: «اگر در هر دو اتاق دختر باشد چطور؟»
	پادشاه گفت: «در این صورت خوشبختی به شما روی می‌آورد.»
	زندانی پرسید: «خب فرض کنیم در یک اتاق ببر و در اتاق دیگر دختر باشد. در آن صورت چه خواهد شد؟»
	پادشاه گفت: «در آن صورت انتخاب شما خیلی اهمیت خواهد داشت.»
	زندانی پرسید: «از کجا بدانم کدام اتاق را باید انتخاب کنم؟»
	پادشاه به علامت‌هایی که بر در اتاق‌ها نصب شده بود اشاره کرد و گفت: «با کمک این علامت‌ها!»
	\begin{center}
		\begin{tikzpicture}
			\draw (0,0) rectangle (4,2.5);
			\node[text width=3cm, align=center] at (2,1.25) {\RL{
		در ۱: در این اتاق یک دختر است و در اتاق دیگر ببر.	
			}};
			\draw (4.5,0) rectangle (8.5,2.5);
			\node[text width=3cm, align=center] at (6.5,1.25) {\RL{
		در ۲: در یکی از این دو اتاق دختر و دیگری ببر است.
			}};
		\end{tikzpicture}
	\end{center}
	زندانی پرسید: «آیا این علامت‌ها درستند؟»
	پادشاه پاسخ داد: «یکی از آن‌ها درست است و دیگری نادرست.»
	اگر شما به جای زندانی بودید کدام در را باز می‌کردید؟ (البته اگر دختر را به ببر ترجیح می‌دهید!)
\end{problem}

\begin{problem} \label{loprpirates}
	به عنوان پاداش نجات دختر پادشاه از دست دزدان دریایی، پادشاه به شما فرصت داده است که گنجی را که درون یکی از سه صندوقچه پنهان شده است، برنده شوید. دو صندوقچه خالی هستند. برای برنده شدن، باید صندوقچه‌ی درست را انتخاب کنید. روی هر کدام از صندوقچه‌ی شماره‌ی ۱ و ۲ این عبارت نوشته شده: «این صندوقچه خالی است.» و روی صندوقچه‌ی شماره‌ی ۳ این عبارت نوشته شده: «گنج در صندوقچه‌ی شماره‌ی ۲ است.» ملکه که هرگز دروغ نمی‌گوید، به شما می‌گوید که تنها یکی از این نوشته‌ها درست است. کدام صندوقچه را باید انتخاب کنید تا برنده شوید؟
\end{problem}

\begin{problem} \label{loprkraig}
	به کارآگاه کریگ مأموریت داده شده تا تیمارستانی را بررسی کند. ساکنان تیمارستان را فقط پزشک و بیمار تشکیل داده است. هر یک از ساکنان - پزشک یا بیمار - یا عاقل است یا دیوانه. عاقل‌ها هر گزاره‌ی درست را درست می‌پندارند و هر گزاره‌ی نادرست را نادرست. دیوانه‌ها کاملاً برعکس هستند. در ضمن، همه‌ی ساکنان تیمارستان راستگو هستند. \\
	نخستین نکته‌ای که کریگ کشف کرد این است که ساکنان تیمارستان کمیته‌های گوناگون تشکیل داده‌اند. در یک کمیته ممکن است هم پزشک و هم بیمار و همچنین هم عاقل و هم دیوانه عضویت داشته باشند. کارگاه به واقعیت‌های دیگری نیز پی برد:
	\begin{enumerate}
		\item 
		همه‌ی بیماران یک کمیته تشکیل داده‌اند.
		\item 
		همه‌ی پزشکان یک کمیته تشکیل داده‌اند.
		\item 
		هر کس در تیمارستان دوستان بسیار دارد که یکی از آن‌ها بهترین دوستش است. همچنین هر کس دشمنان بسیار دارد که یکی از آن‌ها بدترین دشمنش است.
		\item 
		به‌ازای هر کمیته‌ای چون \lr{A}، همه‌ی کسانی که بهترین دوستشان در کمیته‌ی A است، یک کمیته تشکیل داده‌اند. همچنین همه‌ی کسانی که بدترین دشمنشان در کمیته‌ی A است، یک کمیته تشکیل داده‌اند.
		\item 
		به‌ازای هر دو کمیته‌ای چون کمیته‌ی B و کمیته‌ی \lr{C}، دست‌کم شخصی مانند D وجود دارد که بهترین دوستش باور دارد که او در کمیته‌ی B و بدترین دشمنش باور دارد که او در کمیته‌ی C است.
	\end{enumerate}
	کریگ با جمع‌بندی این واقعیت‌ها ثابت کرد که یا یکی از پزشکان دیوانه است یا یکی از بیماران عاقل. او چگونه به این نتیجه رسید؟
\end{problem}

\section{منطق گزاره‌ها}
هر منطق باید موارد زیر را داشته باشد:
\begin{enumerate}
	\item 
	نحوه‌ی نوشتن جملات و عبارات در آن منطق
	\item 
	نحوه‌ای برای نسبت دادن معنا به آن جملات و تخصیص ارزش (درستی یا نادرستی) به آن‌ها
	\item 
	نحوه‌ای برای استدلال کردن در آن منطق
	\item 
	اطمینان از این‌که با شروع از جملات دارای ارزش درست، و با به‌کارگیری نحوه‌ی صحیح استدلال، قطعاً به جملات درست خواهیم رسید.
\end{enumerate}

\subsection{گزاره‌نویسی}

\textbf{گزاره‌ی اتمی:}
یک جمله‌ی خبری ساده است که برای آن یک ارزش درست یا غلط قابل تصور است.
\begin{example}
	جمله‌های زیر گزاره‌ی اتمی هستند:
	\vspace{-0.5cm}
	\begin{multicols}{2}
		\begin{itemize}
			\item 
			عدد ۱۳۷ بزرگتر از عدد ۷۳۱ است.
			\item 
			عدد ۵۰۸۶۱۵۴۳۱۴۵۹۸۴۷ عددی اول است.
		\end{itemize}
	\end{multicols}
\end{example}

\begin{example}
	جمله‌های زیر گزاره‌ی اتمی نیستند:
	\vspace{-0.5cm}
	\begin{multicols}{2}
		\begin{itemize}
			\item 
			عدد طبیعی $n$ عددی زوج است.
			\item 
			آیا مجموع زاویه‌های داخلی مثلث،
			$\ang{180}$
			است؟
		\end{itemize}
	\end{multicols}
\end{example}
\smallskip
\textbf{نحوه‌ی جمله‌سازی در یک منطق گزاره‌ها:}
فرض کنید یک مجموعه‌ی متشکل از گزاره‌های اتمی برای یک منطق گزاره‌ها، ثابت در نظر گرفته شده باشد. مجموعه‌ی متشکل از از همه‌ی گزاره‌های آن منطق گزاره‌ها، با استفاده از قوانین زیر ساخته می‌شوند و به هر عنصر در مجموعه‌ی گزاره‌ها، یک گزاره گفته می‌شود.
\begin{enumerate}
	\item
	هر گزاره‌ی اتمی یک گزاره است. $\top$
	و
	$\bot$
	نیز گزاره هستند.
	\item
	اگر $P$ یک گزاره باشد، آنگاه
	$\neg P$
	نیز یک گزاره است که آن را نقیض $P$ می‌خوانیم.
	\item 
	اگر $P$ و $Q$ دو گزاره باشند، آنگاه عبارت
	$P \wedge Q$
	نیز یک گزاره است که آن را عطف دو گزاره‌ی $P$ و $Q$ می‌نامیم.
	\item 
	اگر $P$ و $Q$ دو گزاره باشند، آنگاه عبارت
	$P \vee Q$
	نیز یک گزاره است که آن را فصل دو گزاره‌ی $P$ و $Q$ می‌نامیم.
	\item 
	اگر $P$ و $Q$ دو گزاره باشند، آنگاه عبارت
	$P \rightarrow Q$
	نیز یک گزاره است که آن را به صورت $P$ آنگاه $Q$ می‌خوانیم.
	\item 
	اگر $P$ و $Q$ دو گزاره باشند، آنگاه عبارت
	$P \leftrightarrow Q$
	نیز یک گزاره است که آن را به صورت $P$ اگر و تنها اگر $Q$ می‌خوانیم.
\end{enumerate}
علائم کمکی استفاده‌شده برای ساختن گزاره‌های بالا، یعنی
$\neg, \wedge, \vee, \rightarrow, \leftrightarrow$
را رابط‌های منطقی می‌نامیم. همچنین، یک گزاره در یک منطق گزاره‌ها، دنباله‌ای متناهی از علائم است.

\begin{example}
	فرض کنید
	$p$، $q$ و $r$
	گزاره‌هایی اتمی در منطق گزاره‌های ما باشند. عبارت‌های زیر گزاره‌ای در منطق گزاره‌های موردنظر ما هستند.
	\vspace{-0.5cm}
	\begin{multicols}{3}
		\begin{itemize}
			\item 
			$\neg p \rightarrow q$
			\item 
			$(p \wedge q) \rightarrow (q \vee r)$
			\item 
			$p \leftrightarrow \neg (r \rightarrow (p \wedge \top))$
		\end{itemize}
	\end{multicols}
\end{example}

\begin{example}
	فرض کنید
	$p$، $q$ و $r$
	گزاره‌هایی اتمی در منطق گزاره‌های ما باشند. عبارت‌های زیر گزاره‌ای در منطق گزاره‌های موردنظر ما نیستند.
	\vspace{-0.5cm}
	\begin{multicols}{2}
		\begin{itemize}
			\item 
			$(p \rightarrow q)\neg$
			\item 
			$p \divideontimes q$
			\item 
			$\forall a \exists b \quad r$
			\item 
			$(((p \vee q) \vee r) \vee \cdots)$
		\end{itemize}
	\end{multicols}
\end{example}

\subsection{معناشناسی}
معناشناسی یک گزاره در منطق گزاره‌ها یعنی رسم جدول ارزش برای آن. T یعنی درست و F یعنی نادرست.

یک گزاره‌ی اتمی $p$ با استفاده از جدول زیر ارزش‌گذاری می‌شود:
\vspace{-0.1cm}
\begin{table}[H]
	\centering
	\begin{tabular}{c}
		$p$ \\
		\hline
		T \\
		F
	\end{tabular}
\end{table}
دو گزاره‌ی
$\top$
و
$\bot$
به صورت زیر ارزش‌گذاری می‌شوند:
\vspace{-0.1cm}
\begin{table}[H]
	\centering
	\begin{tabular}{c}
		$\top$ \\
		\hline
		T
	\end{tabular}
	\qquad
	\begin{tabular}{c}
		$\bot$ \\
		\hline
		F
	\end{tabular}
\end{table}

اگر $P$ یک گزاره باشد، ارزش گزاره‌ی
$\neg P$
به صورت زیر زیر تعریف می‌شود:
\vspace{-0.1cm}
\begin{table}[H]
	\centering
	\begin{latin}
		\begin{tabular}{c|c}
			$P$ & $\neg P$ \\
			\hline
			T & F \\
			F & T
		\end{tabular}
	\end{latin}
\end{table}

اگر $P$ و $Q$ دو گزاره باشند، ارزش‌گذاری گزاره‌ی
$P \wedge Q$
با استفاده از جدول زیر تعریف می‌شود:
\vspace{-0.1cm}
\begin{table}[H]
	\centering
	\begin{latin}
		\begin{tabular}{c|c|c}
			$P$ & $Q$ & $P \wedge Q$ \\
			\hline
			T & T & T \\
			T & F & F \\
			F & T & F \\
			F & F & F
		\end{tabular}
	\end{latin}
\end{table}

اگر $P$ و $Q$ دو گزاره باشند، ارزش‌گذاری گزاره‌ی
$P \vee Q$
با استفاده از جدول زیر تعریف می‌شود:
\vspace{-0.1cm}
\begin{table}[H]
	\centering
	\begin{latin}
		\begin{tabular}{c|c|c}
			$P$ & $Q$ & $P \vee Q$ \\
			\hline
			T & T & T \\
			T & F & T \\
			F & T & T \\
			F & F & F
		\end{tabular}
	\end{latin}
\end{table}

اگر $P$ و $Q$ دو گزاره باشند، ارزش‌گذاری گزاره‌ی
$P \rightarrow Q$
با استفاده از جدول زیر تعریف می‌شود:
\vspace{-0.1cm}
\begin{table}[H]
	\centering
	\begin{latin}
		\begin{tabular}{c|c|c}
			$P$ & $Q$ & $P \rightarrow Q$ \\
			\hline
			T & T & T \\
			T & F & F \\
			F & T & T \\
			F & F & T
		\end{tabular}
	\end{latin}
\end{table}
در گزاره‌ی
$P \rightarrow Q$
به $P$ مقدم و به $Q$ تالی گفته می‌شود. هرگاه ارزش مقدم نادرست (F) باشد، می‌گوییم گزاره‌ی
$P \rightarrow Q$
به انتفای مقدم درست (T) است.

اگر $P$ و $Q$ دو گزاره باشند، ارزش‌گذاری گزاره‌ی
$P \leftrightarrow Q$
با استفاده از جدول زیر تعریف می‌شود:
\vspace{-0.1cm}
\begin{table}[H]
	\centering
	\begin{latin}
		\begin{tabular}{c|c|c}
			$P$ & $Q$ & $P \leftrightarrow Q$ \\
			\hline
			T & T & T \\
			T & F & F \\
			F & T & F \\
			F & F & T
		\end{tabular}
	\end{latin}
\end{table}

\begin{example}
	نقیض گزاره‌های زیر را بنویسید.
	\vspace{-0.5cm}
	\begin{multicols}{2}
		\begin{enumerate}
			\item 
			عدد ۲ کوچک‌تر از عدد ۳ نیست.
			\item 
			عدد ۴ زوج است.
		\end{enumerate}
	\end{multicols}
	\noindent \textbf{پاسخ:}
	\vspace{-0.5cm}
	\begin{multicols}{2}
		\begin{enumerate}
			\item 
			عدد ۲ بزرگ‌تر یا مساوی عدد ۳ است.
			\item 
			عدد ۴ فرد است. (عدد ۴ زوج نیست.)
		\end{enumerate}
	\end{multicols}
	\vspace{0.001cm}
\end{example}

\begin{example}
	جدول ارزش گزاره‌ی
	$(p \vee q) \wedge \neg p$
	را رسم کنید.
	\begin{table}[H]
		\centering
		\begin{latin}
			\begin{tabular}{c|c|c|c|c}
				$p$ & $q$ & $p \vee q$ & $\neg p$ & $(p \vee q) \wedge \neg p$ \\
				\hline
				T & T & T & F & F \\
				T & F & T & F & F \\
				F & T & T & T & T \\
				F & F & F & T & F
			\end{tabular}
		\end{latin}
	\end{table}
\end{example}

\begin{example}
	جدول ارزش گزاره‌ی
	$(p \rightarrow q) \rightarrow \neg(q \rightarrow p)$
	را رسم کنید.
	\begin{table}[H]
		\centering
		\begin{latin}
			\begin{tabular}{c|c|c|c|c|c}
				$p$ & $q$ & $p \rightarrow q$ & $q \rightarrow p$ & $\neg(q \rightarrow p)$ & $(p \rightarrow q) \rightarrow \neg(q \rightarrow p)$ \\
				\hline
				T & T & T & T & F & F \\
				T & F & F & T & F & T \\
				F & T & T & F & T & T \\
				F & F & T & T & F & F \\
			\end{tabular}
		\end{latin}
	\end{table}
\end{example}

\begin{example}[پیش‌نیاز: احتمال]
	فرض کنید $p$ و $q$ گزاره‌های اتمی باشند. اگر بدانیم
	$p \rightarrow q$
	درست است، آنگاه احتمال درستی
	$q \leftrightarrow p$
	چقدر است؟
	\\ \textbf{پاسخ:}
	\begin{table}[H]
		\centering
		\begin{latin}
			\begin{tabular}{c|c|c|c}
				$p$ & $q$ & $p \rightarrow q$ & $q \leftrightarrow p$ \\
				\hline
				T & T & T & T \\
				T & F & F & F \\
				F & T & T & F \\
				F & F & T & T
			\end{tabular}
		\end{latin}
	\end{table}
	\noindent
	حالت‌های ممکن: سطرهای ۱، ۳ و ۴، و حالت‌های مطلوب: سطرهای ۱ و ۴. احتمال درستی گزاره‌ی
	$q \leftrightarrow p$
	برابر است با تعداد حالت‌های مطلوب تقسیم بر تعداد حالت‌های ممکن. بنابراین، این گزاره به احتمال
	$\dfrac{2}{3}$
	درست است.
\end{example}

\subsection{گزاره‌های معادل}
دو گزاره‌ی $P$ و $Q$ در منطق گزاره‌ها معادل یا هم‌ارز هستند، هرگاه وقتی جدول ارزش آن دو برحسب گزاره‌های اتمی به کار رفته در آن‌ها کشیده شود، ستون آخر یکسان درآید؛ یعنی زمانی که گزاره‌های اتمی استفاده‌شده در این دو گزاره، ارزش یکسان داشته باشند، ارزش این دو گزاره یکسان شود. در این صورت می‌نویسیم:
$$P \iff Q$$
عبارت بالا یک گزاره در منطق گزاره‌ها نیست.
\smallskip
\begin{example}
	نشان دهید که
	$p \leftrightarrow q \iff (p \rightarrow q) \wedge (q \rightarrow p)$.
	\begin{table}[H]
		\centering
		\begin{latin}
			\begin{tabular}{c|c|c|c|c|c}
				$p$ & $q$ & $p \rightarrow q$ & $q \rightarrow p$ & $(p \rightarrow q) \wedge (q \rightarrow p)$ & $p \leftrightarrow q$ \\
				\hline
				T & T & T & T & T & T \\
				T & F & F & T & F & F \\
				F & T & T & F & F & F \\
				F & F & T & T & T & T
			\end{tabular}
		\end{latin}
	\end{table}
\end{example}

\begin{example}
	نشان دهید که گزاره‌ی
	$p \rightarrow q$
	با گزاره‌ی عکس نقیض خود معادل است. \\
	\textbf{پاسخ:}
	باید نشان دهیم که
	$p \rightarrow q \iff \neg q \rightarrow \neg p$.
	\begin{table}[H]
		\centering
		\begin{latin}
			\begin{tabular}{c|c|c|c|c|c}
				$p$ & $q$ & $\neg q$ & $\neg p$ & $\neg q \rightarrow \neg p$ & $p \rightarrow q$ \\
				\hline
				T & T & F & F & T & T \\
				T & F & T & F & F & F \\
				F & T & F & T & T & T \\
				F & F & T & T & T & T
			\end{tabular}
		\end{latin}
	\end{table}
\end{example}
\smallskip
\textbf{استلزام منطقی:}
هرگاه گزاره‌ی
$P \rightarrow Q$
همواره دارای ارزش درست (T) باشد، می‌نویسیم:
$$P \implies Q$$
و می‌خوانیم: «$P$ مستلزم $Q$ است.» یا «$P$ شرط کافی برای $Q$ است.» یا «$Q$ شرط لازم برای $P$ است.» \\
عبارت بالا نیز یک گزاره در منطق گزاره‌ها نیست.

\begin{theorem}
	$P \iff Q$
	اگر و تنها اگر
	$P \implies Q$ و $Q \implies P$.
\end{theorem}

\begin{theorem} \label{lotheq}
	در منطق گزاره‌ها چنین است:
	\vspace{-0.3cm}
	\begin{multicols}{2}
		\begin{enumerate}
			\item 
			$p \rightarrow q \iff \neg p \vee q$
			\item 
			$p \vee (q \vee r) \iff (p \vee q) \vee r$
			\item 
			$p \wedge (q \wedge r) \iff (p \wedge q) \wedge r$
			\item 
			$p \vee q \iff q \vee p$
			\item 
			$p \wedge q \iff q \wedge p$
			\item 
			$p \wedge (q \vee r) \iff (p \wedge q) \vee (p \wedge r)$
			\item 
			$p \vee (q \wedge r) \iff (p \vee q) \wedge (p \vee r)$
			\item 
			$p \wedge \top \iff p$
			\item 
			$p \vee \top \iff \top$
			\item 
			$p \wedge \bot \iff \bot$
			\item 
			$p \vee \bot \iff p$
			\item 
			$p \wedge p \iff p$
			\item 
			$p \vee p \iff p$
			\item 
			$p \wedge (p \vee q) \iff p$
			\item 
			$p \vee (p \wedge q) \iff p$
			\item 
			$p \wedge \neg p \iff \bot$
			\item 
			$p \vee \neg p \iff \top$
			\item 
			$\neg (\neg p) \iff p$
			\item 
			$\neg (p \wedge q) \iff \neg p \vee \neg q$
			\item 
			$\neg (p \vee q) \iff \neg p \wedge \neg q$
			\item 
			$\neg (p \rightarrow q) \iff p \wedge \neg q$
			\item 
			$\neg (p \leftrightarrow q) \iff \neg p \leftrightarrow q \iff p \leftrightarrow \neg q$
		\end{enumerate}
	\end{multicols}
	\vspace{0.001cm}
\end{theorem}
\smallskip
برای اثبات معادل بودن دو گزاره، دو راه وجود دارد. یک روش کشیدن جدول ارزش است و روش دیگر استفاده از قوانینی محدود است، برای مثال استفاده از قوانینی که در قضیه‌ی
\ref{lotheq}
بیان کرده‌ایم. به روش اول، استلزام و به روش دوم، استنتاج می‌گویند.

\begin{example}
	نشان دهید که
	$p \rightarrow q \iff (p \wedge \neg q) \rightarrow \bot$.
	\\ \textbf{پاسخ:}
	\begin{align*}
		(p \wedge \neg q) \rightarrow \bot &\iff \neg(p \wedge \neg q) \vee \bot &\text{بنا به قانون ۱} \\
		&\iff \neg(p \wedge \neg q) &\text{بنا به قانون ۱۱} \\
		&\iff \neg p \vee q &\text{بنا به قانون‌های ۱۸ و ۱۹} \\
		&\iff p \rightarrow q &\text{بنا به قانون ۱}
	\end{align*}
\end{example}
\begin{example}
	اگر
	$(p \vee \neg q) \rightarrow \neg r$ و $(\neg p \wedge q) \rightarrow \neg r$
	گزاره‌هایی درست باشند، آنگاه درباره‌ی ارزش گزاره‌ی
	$(r \vee q) \rightarrow (r \wedge p)$
	چه می‌توان گفت؟
	\\ \textbf{پاسخ:}
	می‌دانیم که
	$p \vee \neg q \iff \neg (\neg p \wedge q)$.
	داریم:
	\begin{align*}
		((\neg p \wedge q) \rightarrow \neg r) \wedge (p \vee \neg q) \rightarrow \neg r
		&\iff ((p \vee \neg q) \vee \neg r) \wedge ((\neg p \wedge q) \vee \neg r) \\
		&\iff \neg r \vee ((p \vee \neg q) \wedge (\neg p \wedge q)) \\
		&\iff \neg r \vee \bot \\
		&\iff \neg r
	\end{align*}
	چون گزاره‌های مفروض سؤال درست هستند، عطف آن‌ها نیز درست است، پس
	$\neg r \iff \top$.
	یعنی $r$ دارای ارزش نادرست (F) است. بنابراین:
	\begin{align*}
		(r \vee q) \rightarrow (r \wedge p)
		&\iff (\bot \vee q) \rightarrow (\bot \wedge p) \\
		&\iff q \rightarrow \bot
	\end{align*}
	در نتیجه، ارزش گزاره‌ی خواسته‌شده به ارزش $q$ بستگی دارد. اگر $q$ درست باشد، آنگاه آن گزاره نادرست است و اگر $q$ نادرست باشد، آنگاه آن گزاره درست است. به بیان دیگر:
	$$(r \vee q) \rightarrow (r \wedge p) \iff \neg q$$
\end{example}

\section{منطق مرتبه‌ی اول}
در این قسمت، منطق مرتبه‌ی اول را بسیار ناقص معرفی می‌کنیم. اگر علاقه‌مند به یادگیری دقیق هستید می‌توانید به کتاب‌ها و جزوه‌های معتبر منطق ریاضی مراجعه کنید. همچنین، پیش‌نیاز این قسمت، سایر فصل‌ها هستند، بنابراین می‌توانید بعداً به این قسمت بازگردید و یاد بگیرید.

\subsection{گزاره‌نما}
هر جمله‌ی خبری که شامل یک یا چند متغیر است و با جای‌گذاری مقادیری به جای متغیر به یک گزاره تبدیل شود، گزاره‌نما نامیده می‌شود. گزاره‌نماها را برحسب تعداد متغیر به کار رفته در آن‌ها، یک‌متغیره، دومتغیره و ... می‌نامیم.

\begin{example}
	جملات خبری زیر گزاره‌نما هستند.
	\vspace{-0.5cm}
	\begin{multicols}{2}
		\begin{enumerate}
			\item 
			$a$
			عددی فرد است.
			\item 
			دو برابر $x$ مساوی ۱۰ است.
		\end{enumerate}
	\end{multicols}
	اگر در جمله‌ی «۱» قرار دهیم
	$a=1$
	در این صورت ارزش گزاره‌ی حاصل درست (T) می‌شود. \\
	اگر در جمله‌ی «۲» قرار دهیم
	$x=4$
	در این صورت ارزش گزاره‌ی حاصل نادرست (F) می‌شود.
\end{example}
\smallskip
در هر گزاره‌نما به مجموعه‌ی مقادیری که می‌توان آن‌ها را به‌جای متغیرهای آن قرار داد تا گزاره‌نما به گزاره تبدیل شود، دامنه‌ی متغیر گزاره‌نما می‌گویند و آن را با حرف $D$ نمایش می‌دهند. در هر گزاره‌نما، به مجموعه‌ی عضوهایی از دامنه‌ی متغیر که به‌ازای آن‌ها، گزاره‌نما تبدیل به گزاره‌ای با ارزش درست (T) شود، مجموعه‌ی جواب گزاره‌نما می‌گویند و آن را با حرف $S$ نمایش می‌دهند و همواره داریم:
$S \subseteq D$.

\begin{example}
	گزاره‌نمای «$m$ عددی زوج است» را در نظر بگیرید. دامنه‌ی متغیر آن مجموعه‌ی اعداد صحیح مثبت
	($\mathbb{Z}^+$)
	است و مجموعه‌ی جواب آن
	$m=2k$
	که $k$ عددی صحیح و مثبت است. 
\end{example}

\subsection{سورها}
عبارت‌هایی مانند «همه»، «هر»، «به‌ازای هر»، «به‌ازای بعضی مقدارها»، «وجود دارد» به سور معروف هستند. برای بیان عبارت‌ها با نمادهای ریاضی به‌جای «هر» از نماد
$\forall$
و به‌جای «وجود دارد» از نماد
$\exists$
استفاده می‌کنیم که به ترتیب، سور عمومی و سور وجودی نامیده می‌شوند.
\begin{example}
	فرض کنید در جهان اعداد حقیقی، برقراری
	$A(x,y)$
	یعنی $x$ کوچک‌تر از $y$ است. می‌خواهیم جمله‌ی «یک عدد هست که از همه‌ی اعداد بزرگ‌تر است.» را به زبان ریاضی بنویسیم.
	$$\exists x \ \forall y \ A(y,x)$$
\end{example}
\begin{example}
	جمله‌ی
	$\forall x \in \mathbb{R} \ \exists n \in \mathbb{N} \ \ x \leq n$
	را به زبان فارسی بنویسید.
	\\ \textbf{پاسخ:}
	به‌ازای هر $x$ متعلق به مجموعه‌ی اعداد حقیقی، یک $n$ متعلق به مجموعه‌ی اعداد طبیعی وجود دارد که $x$ کوچک‌تر یا مساوی $n$ است.
\end{example}
\smallskip
گزاره‌نمایی که با سور عمومی بیان می‌شود، فقط وقتی درست است که همیشه و برای همه‌ی اعضا درست باشند؛ یا به عبارت دیگر، دامنه‌ی متغیر آن مساوی با مجموعه‌ی جوابش باشد. گزاره‌نمایی که با سور وجودی بیان می‌شود، فقط وقتی درست است که حداقل دارای یک جواب درست باشد؛ یا به عبارت دیگر مجموعه‌ی جواب آن تهی نباشد.
\begin{theorem}
	در منطق مرتبه‌ی اول چنین است:
	\begin{enumerate}
		\item 
		$\exists x \ p(x) \iff \exists y \ p(y)$
		\item 
		$\forall x \ p(x) \iff \forall y \ p(y)$
		\item 
		$\neg(\forall x \ p(x)) \iff \exists x \ (\neg p(x))$
		\item 
		$\neg (\exists x \ p(x)) \iff \forall x \ (\neg p(x))$
		\item 
		$\forall x \ (p(x) \wedge q(x)) \iff ((\forall x \ p(x)) \wedge (\forall x \ q(x)))$
		\item 
		$\exists x \ (p(x) \vee q(x)) \iff ((\exists x \ p(x)) \vee (\exists x \ q(x)))$
	\end{enumerate}
\end{theorem}
\begin{example}
	نقیض جمله‌های زیر را بنویسید.
	\vspace{-0.5cm}
	\begin{multicols}{2}
		\begin{enumerate}
			\item 
			هر عدد زوج بر ۲ بخش‌پذیر است.
			\item 
			بعضی مثلث‌ها متساوی‌الساقین هستند.
			\item 
			هیچ مجموعه‌ای متعلق به خودش نیست.
			\item 
			عدد اولی که مضرب ۴ نباشد، وجود دارد.
		\end{enumerate}
	\end{multicols}
	\noindent \textbf{پاسخ:}
	\vspace{-0.5cm}
	\begin{multicols}{2}
		\begin{enumerate}
			\item 
			عدد زوجی وجود دارد که بر ۲ بخش‌پذیر نیست.
			\item 
			هیچ مثلثی متساوی‌الساقین نیست.
			\item 
			مجموعه‌ای وجود دارد که متعلق به خودش است.
			\item 
			همه‌ی عددهای اول مضرب ۴ هستند.
		\end{enumerate}
	\end{multicols}
	\vspace{0.1pt}
\end{example}
\vspace{0.5cm}
\section{مسئله‌ها}
\begin{problem} \label{lorjnzmg}
	ساکنان یک جزیره دو دسته‌اند، یا «نجیب‌زاده»‌اند، کسانی که همیشه راست می‌گویند، یا «متقلب»اند، کسانی که همیشه دروغ می‌گویند. چند تن از جزیره‌نشینان در همایشی یکدیگر را دیدند و هر یک از آن‌ها به دیگری گفت: «شما همگی متقلب‌اید.» چند نجیب‌زاده ممکن است در این همایش شرکت کرده باشند؟
\end{problem}

\begin{problem} \label{lopr4cab45}
	چهار کارت روی میزی وجود دارند که روی طرف قابل رؤیتشان نمادهای \lr{A}، \lr{‌B}، ۴ و ۵ نوشته شده است. دست‌کم چند کارت (و کدام کارت) را باید برگردانیم تا بفهمیم گزاره‌ی زیر درست است یا نه؟ \\
	«اگر روی یک طرف کارتی عددی زوج نوشته شده باشد، آنگاه روی طرف دیگرش یک حرف صدادار نوشته شده است.»
\end{problem}

\begin{problem} \label{loprtvmth}
	فرض کنید دو گزاره‌ی زیر درست باشند:
	\begin{enumerate}
		\item 
		در میان افرادی که تلویزیون دارند، کسانی هستند که ریاضی‌دان نیستند.
		\item 
		کسی که ریاضی‌دان نیست و هر روز در استخر شنا می‌کند، تلویزیون ندارد.
	\end{enumerate}
	آیا می‌توانیم ادعا کنیم که «این‌طور نیست که همه‌ی افرادی که تلویزیون دارند، هر روز شنا می‌کنند»؟
\end{problem}

\begin{problem} \label{loprdpesffs}
	در دبیرستان پسرانه‌ی ابن‌سینا تعدادی دانش‌آموز تحصیل می‌کنند. اگر دو دانش‌آموز از این مدرسه دوست باشند، پدرهای این دو دانش‌آموز نیز با هم دوست هستند (دوستی یک رابطه‌ی دوطرفه است و هیچکس با خودش دوست نیست). کدام گزینه قطعاً درست است؟
	\begin{enumerate}
		\item 
		اگر حامد (یکی از دانش‌آموزان مدرسه) $k$ دوست در دبیرستان داشته باشد، آنگاه پدر حامد حداقل $k$ دوست در میان پدرهای دانش‌آموزان دبیرستان دارد.
		\item 
		اگر حامد $k$ دوست در دبیرستان داشته باشد، آنگاه پدر حامد حداکثر $k$ دوست در میان پدرهای دانش‌آموزان دبیرستان دارد.
		\item 
		اگر $k$ دانش‌آموز از دبیرستان موجود باشند که هیچ دوتایی از آن‌ها دوست نباشند، آنگاه در میان پدرهای دانش‌آموزان دبیرستان، $k$ نفر وجود دارند که هیچ دوتایی دوست نیستند.
		\item 
		اگر در میان پدرهای دانش‌آموزان دبیرستان، $k$ نفر موجود باشند که هیچ دوتایی دوست نباشند، آنگاه $k$ دانش‌آموز وجود دارند که هیچ دوتایی دوست نیستند.
		\item 
		اگر حسام برادر کوچک‌تر حامد باشد، آنگاه حسام و حامد با هم دوست هستند (حسام و حامد هر دو در دبیرستان ابن‌سینا تحصیل می‌کنند).
	\end{enumerate}
\end{problem}

\begin{problem} \label{loprrttp}
	جدول ارزش‌های گزاره‌های زیر را رسم کنید.
	\vspace{-5mm}
	\begin{multicols}{2}
		\begin{enumerate}
			\item 
			$(p \vee q) \rightarrow q$
			\item 
			$\neg p \wedge (p \leftrightarrow q)$
		\end{enumerate}
	\end{multicols}
	\vspace{1pt}
\end{problem}

\begin{problem} \label{loprpttp}
	با استفاده از جدول ارزش‌ها نشان دهید که:
	\begin{enumerate}
		\item 
		$\neg(p \leftrightarrow q) \iff \neg p \leftrightarrow q$
		\item 
		$p \rightarrow (q \rightarrow r) \iff (p \wedge q) \rightarrow r$
	\end{enumerate}
\end{problem}

\begin{problem} \label{loprestentaj}
	با استفاده از گزاره‌های معادل (قضیه‌ی \ref{lotheq}) نشان دهید که:
	\vspace{-5mm}
	\begin{multicols}{2}
		\begin{enumerate}
			\item 
			$p \rightarrow \bot \iff \neg p$
			\item 
			$(p \rightarrow q) \wedge q \iff p \wedge q$
		\end{enumerate}
	\end{multicols}
	\vspace{1pt}
\end{problem}

\begin{problem} \label{loprpqreqx}
	گزاره‌ای شامل گزاره‌های اتمی $p$، $q$ و $r$ بنویسید که هم‌ارز منطقی گزاره‌ی $X$ باشد.
	\vspace{-2mm}
	\begin{table}[H]
		\centering
		\begin{latin}
			\begin{tabular}{c|c|c|c}
				$p$ & $q$ & $r$ & $X$ \\
				\hline
				T & T & T & F \\
				T & T & F & F \\
				T & F & T & T \\
				T & F & F & F \\
				F & T & T & F \\
				F & T & F & T \\
				F & F & T & T \\
				F & F & F & F
			\end{tabular}
		\end{latin}
	\end{table}
\end{problem}

\begin{problem}[پیش‌نیاز: احتمال] \label{loprttpr}
	ارزش گزاره‌ی
	$p \rightarrow (q \vee r)$
	درست است. احتمال این‌که ارزش گزاره‌ی $r$ نادرست باشد، چقدر است؟
\end{problem}

\begin{problem} [پیش‌نیاز: حسابان] \label{loprcafcd}
	دو جمله‌ی زیر را در نظر بگیرید:
	\begin{enumerate}
		\item 
		اگر تابع $f$ در نقطه‌ی $a$ پیوسته باشد، آنگاه تابع $f$ یک تابع مطلوب است.
		\item 
		اگر تابع $f$ در نقطه‌ی $a$ مشتق‌پذیر باشد، آنگاه تابع $f$ یک تابع مطلوب است.
	\end{enumerate}
	کدام جمله‌ی بالا از دیگری نتیجه می‌شود؟
\end{problem}
